\documentclass[11pt,a4paper,openany]{book}

\usepackage[utf8]{inputenc}
\usepackage[T1]{fontenc}
\usepackage{amsmath,amssymb,amsthm,mathrsfs}
\usepackage{geometry}
\geometry{margin=1in,top=1.2in,bottom=1.2in}
\usepackage{hyperref}
\hypersetup{colorlinks=true,linkcolor=blue,citecolor=blue,urlcolor=blue}
\usepackage{enumitem}
\usepackage{booktabs}
\usepackage{array}
\usepackage{tabularx}
\usepackage{mdframed}
\usepackage{xcolor}
\usepackage{tocloft}
\usepackage{fancyhdr}
\usepackage{titlesec}
\usepackage{graphicx}
\setcounter{tocdepth}{2}
\setlength{\parindent}{0pt}
\setlength{\parskip}{0.55em}

\pagestyle{fancy}
\fancyhf{}
\fancyhead[LE]{\small\leftmark}
\fancyhead[RO]{\small\rightmark}
\fancyfoot[C]{\thepage}
\renewcommand{\headrulewidth}{0.3pt}

% Theorem environments
\newtheorem{theorem}{Theorem}[chapter]
\newtheorem{proposition}[theorem]{Proposition}
\newtheorem{lemma}[theorem]{Lemma}
\newtheorem{corollary}[theorem]{Corollary}
\theoremstyle{definition}
\newtheorem{definition}[theorem]{Definition}
\newtheorem{remark}[theorem]{Remark}
\newtheorem{example}[theorem]{Example}

% Boxes
\newmdenv[linecolor=black,backgroundcolor=gray!8,linewidth=0.7pt,
  innertopmargin=7pt,innerbottommargin=7pt,
  innerleftmargin=9pt,innerrightmargin=9pt,
  skipabove=8pt,skipbelow=8pt]{lockedbox}

\newmdenv[linecolor=blue!70,backgroundcolor=blue!5,linewidth=0.7pt,
  innertopmargin=7pt,innerbottommargin=7pt,
  innerleftmargin=9pt,innerrightmargin=9pt,
  skipabove=8pt,skipbelow=8pt]{derivbox}

\newmdenv[linecolor=red!60,backgroundcolor=red!4,linewidth=0.7pt,
  innertopmargin=7pt,innerbottommargin=7pt,
  innerleftmargin=9pt,innerrightmargin=9pt,
  skipabove=8pt,skipbelow=8pt]{resultbox}

% Macros
\newcommand{\R}{\mathbb{R}}
\newcommand{\C}{\mathbb{C}}
\newcommand{\T}{^{\mathsf{T}}}
\newcommand{\norm}[1]{\left\|#1\right\|}
\newcommand{\inner}[2]{\langle #1,\,#2\rangle}
\newcommand{\Fbase}{F_{\mathrm{base}}}
\newcommand{\gX}{g_X}
\newcommand{\half}{\tfrac{1}{2}}
\newcommand{\diag}{\mathrm{diag}}
\newcommand{\tr}{\mathrm{tr}}
\newcommand{\rank}{\mathrm{rank}}
\newcommand{\mfls}{\mathrm{MFLS}}
\newcommand{\adv}{\mathrm{adv}}
\newcommand{\nom}{\mathrm{nom}}
\newcommand{\rob}{\mathrm{rob}}
\newcommand{\so}{\mathfrak{so}}
\newcommand{\fg}{\mathfrak{g}}
\newcommand{\SE}{\mathrm{SE}}
\newcommand{\SO}{\mathrm{SO}}
\newcommand{\SU}{\mathrm{SU}}

\title{
  \Huge\textbf{The Trigonometry of Collapse}\\[0.5em]
  \Large Angular Geometry, Phase Dynamics, and\\
  Rotational Structure in Complex System Instability\\[1em]
  \large A Complete Mathematical Treatment of the\\
  Canonical--BSDT Angular Framework
}
\author{}
\date{}

\begin{document}

\frontmatter
\maketitle

\begin{center}
\vspace{2em}
\begin{minipage}{0.82\textwidth}
\textit{
Most mathematical frameworks for complex system failure monitor scalar magnitudes:
energy, distance, variance, likelihood. This book develops a complementary view:
collapse is fundamentally an angular phenomenon. Systems do not simply grow unstable —
they lose their direction. The gradient and the force cease to point the same way.
Agents synchronise their phases into a collective that the restoring geometry cannot hold.
The critical manifold is not a threshold of magnitude but a surface where curvature
overwhelms gradient. This book derives these statements rigorously, building from
the three fundamental angles of the canonical framework to a complete trigonometric
architecture for early warning, phase prediction, and collapse characterisation.
}
\end{minipage}
\end{center}

\vspace{3em}
\tableofcontents

\mainmatter

%==============================================================
\part{The Three Fundamental Angles}
%==============================================================

\chapter{The Alignment Angle $\theta_t$}

\section{Definition and geometric origin}

The canonical dynamical-geometry engine operates by projecting the nominal force field
$\Fbase$ onto the energy gradient $\gX = 2J\T GS$. The quality of this alignment
determines the engine's effectiveness. The alignment angle quantifies it.

\begin{definition}[Alignment angle]
\[
\cos\theta_t = \frac{\inner{\gX(t)}{\Fbase(t)}}{\norm{\gX(t)}\,\norm{\Fbase(t)}} \in [-1, +1],
\qquad
\theta_t = \arccos(\cos\theta_t) \in [0, \pi].
\]
\end{definition}

\begin{lockedbox}
$\theta_t = 0$: $\Fbase$ points exactly along $\gX$ --- nominal force fully drives instability.\\
$\theta_t = \pi/2$: $\Fbase \perp \gX$ --- force tangential, no first-order energy change.\\
$\theta_t = \pi$: $\Fbase$ anti-parallel to $\gX$ --- force actively opposes instability.
\end{lockedbox}

\section{The energy derivative in terms of $\theta_t$}

\begin{theorem}[Energy derivative via alignment]
\[
\dot E = (1-\gamma)\,\norm{\Fbase}\,\norm{\gX}\cos\theta_t - (1-\gamma)\norm{\gX}^2.
\]
\end{theorem}

\begin{proof}
From the canonical ODE, $\dot E = (1-\gamma)[\inner{\gX}{\Fbase} - \norm{\gX}^2]$.
Substituting $\inner{\gX}{\Fbase} = \norm{\gX}\norm{\Fbase}\cos\theta_t$ gives the result.
\end{proof}

\begin{corollary}[Descent condition in angular form]
\[
\dot E \leq 0 \iff \cos\theta_t \leq \frac{\norm{\gX}}{\norm{\Fbase}} = \frac{\mfls}{\norm{\Fbase}}.
\]
Descent is guaranteed whenever the alignment angle exceeds the critical value
$\theta_t^\star = \arccos(\mfls/\norm{\Fbase})$.
\end{corollary}

\section{Force decomposition via $\theta_t$}

\begin{proposition}[Force decomposition]
The nominal force decomposes into collapse-directed and safe components:
\[
\Fbase = F_\parallel + F_\perp,
\qquad
F_\parallel = \norm{\Fbase}\cos\theta_t\,\hat u,
\qquad
F_\perp = \Fbase - \norm{\Fbase}\cos\theta_t\,\hat u,
\]
where $\hat u = \gX/\norm{\gX}$. Then:
\[
\norm{F_\parallel} = \norm{\Fbase}\cos\theta_t, \quad
\norm{F_\perp} = \norm{\Fbase}\sin\theta_t, \quad
\norm{\Fbase}^2 = \norm{F_\parallel}^2 + \norm{F_\perp}^2.
\]
The canonical ODE damps $F_\parallel$ by factor $(1-\gamma)$ and leaves $F_\perp$ free.
\end{proposition}

\begin{proof}
$\norm{F_\parallel}^2 = \norm{\Fbase}^2\cos^2\theta_t$, $\norm{F_\perp}^2 = \norm{\Fbase}^2 - \norm{F_\parallel}^2 = \norm{\Fbase}^2\sin^2\theta_t$. Pythagorean identity: $\cos^2\theta_t + \sin^2\theta_t = 1$.
\end{proof}

\section{Angular velocity and acceleration}

\begin{theorem}[Angular velocity, exact formula]
\label{thm:angvel}
Let $h(t) = \cos\theta_t$, $p = \inner{\gX}{\Fbase}$, $a = \norm{\gX}$, $b = \norm{\Fbase}$.
Then $\dot\theta_t = -\dot h / \sin\theta_t$ where:
\[
\dot h = \frac{\inner{\dot\gX}{\Fbase} + \inner{\gX}{\dot\Fbase}}{ab}
- \cos\theta_t\!\left(\frac{\inner{\gX}{\dot\gX}}{a^2} + \frac{\inner{\Fbase}{\dot\Fbase}}{b^2}\right).
\]
\end{theorem}

\begin{proof}
Differentiate $h = p/(ab)$ by the quotient rule:
$\dot h = (\dot p \cdot ab - p(\dot a b + a\dot b))/(ab)^2 = \dot p/(ab) - h(\dot a/a + \dot b/b)$.
Then $\dot p = \inner{\dot\gX}{\Fbase} + \inner{\gX}{\dot\Fbase}$,
$\dot a = \inner{\gX}{\dot\gX}/a$, $\dot b = \inner{\Fbase}{\dot\Fbase}/b$.
Substitute to obtain the stated formula. Since $\theta_t = \arccos(h)$,
$\dot\theta_t = -\dot h/\sqrt{1-h^2} = -\dot h/\sin\theta_t$ for $\theta_t \in (0,\pi)$.
\end{proof}

\begin{definition}[Angular acceleration]
\[
\ddot\theta_t = \frac{d^2\theta_t}{dt^2} = \frac{d\dot\theta_t}{dt}.
\]
In discrete time: $\ddot\theta_t \approx (\theta_{t+1} - 2\theta_t + \theta_{t-1})/\Delta t^2$.
\end{definition}

\begin{proposition}[Sign conventions]
$\dot\theta_t > 0$: angle opening, force rotating away from gradient.\\
$\dot\theta_t < 0$: angle closing, force aligning with gradient.\\
$\ddot\theta_t > 0$: decoherence accelerating.\\
$\ddot\theta_t < 0$: decoherence decelerating (angular dynamics stabilising).
\end{proposition}

\section{The alignment angle as a stability indicator}

\begin{theorem}[Safety ratio and alignment]
The safety ratio satisfies:
\[
\mathrm{SafetyRatio}(t) = \frac{\theta\,\mfls_t}{2M_\mathrm{max}\sqrt{E_t}}.
\]
Its time derivative is:
\[
\frac{d}{dt}\mathrm{SafetyRatio} = \frac{\theta}{2M_\mathrm{max}}
\frac{\dot\mfls_t\sqrt{E_t} - \mfls_t\dot E_t/(2\sqrt{E_t})}{E_t}.
\]
The safety ratio falls whenever $\dot\mfls_t/\mfls_t < \dot E_t/(2E_t)$,
i.e.\ gradient grows slower than energy.
\end{theorem}

\begin{proof}
Differentiate $\mathrm{SafetyRatio} = \theta\mfls/(2M_\mathrm{max}\sqrt{E})$
by the quotient rule. The condition $d\,\mathrm{SafetyRatio}/dt < 0$ rearranges to
$2E_t\dot\mfls_t < \mfls_t\dot E_t$, i.e.\ $\dot\mfls_t/\mfls_t < \dot E_t/(2E_t)$.
\end{proof}

\section{Relationship between $\theta_t$ and the admissibility threshold}

From $\Psi^* = \theta\mfls/(4\sqrt{E})$ and the descent condition
$\cos\theta_t \leq \mfls/\norm{\Fbase}$:

\begin{corollary}[Angular admissibility]
The system is admissible ($M_\mathrm{max} < \Psi^*$) if and only if:
\[
\cos\theta_t^\star \leq \frac{2M_\mathrm{max}\sqrt{E}}{\theta\norm{\Fbase}}
\]
at the critical alignment angle. The system breaches admissibility when
the alignment angle falls below this threshold simultaneously with energy rising.
\end{corollary}

%--------------------------------------------------------------
\chapter{The Misalignment Angle $\psi_t$}
%--------------------------------------------------------------

\section{Definition}

The canonical ODE uses the state-space gradient $\tilde G_t = J\T g_t$ (the pullback of the
channel-space gradient through the Jacobian). The original BSDT master operator used the
channel-space gradient $g_t$ directly. These two vectors are generally not parallel.
The angle between them is $\psi_t$.

\begin{definition}[Misalignment angle]
\label{def:psi}
\[
\cos\psi_t = \frac{\inner{\tilde G_t}{g_t}_F}{\norm{\tilde G_t}_F\,\norm{g_t}}
= \frac{R_t}{\norm{\tilde G_t}_F\,\norm{g_t}}
\]
where $\tilde G_t = J_t\T g_t \in \R^{Nd}$ is the pullback gradient,
$g_t = \nabla_S E_\mathrm{BS}(S_t) \in \R^4$ is the channel-space gradient, and
$R_t = g_t\T \mathcal{G}_t g_t = g_t\T J_t J_t\T g_t \geq 0$ is the control coupling sum.
\end{definition}

\section{The Jacobian as an angle-generating map}

The Jacobian $J_t \in \R^{4\times Nd}$ maps channel-space vectors to state-space vectors:
$J_t\T : \R^4 \to \R^{Nd}$.

\begin{proposition}[Singular value decomposition of $J$]
Let $J_t = U\Sigma V\T$ be the thin SVD with $U \in \R^{4\times 4}$, $\Sigma \in \R^{4\times 4}$,
$V \in \R^{Nd\times 4}$. Then:
\[
\tilde G_t = J_t\T g_t = V\Sigma U\T g_t.
\]
The misalignment angle satisfies:
\[
\cos\psi_t = \frac{(U\T g_t)\T\Sigma^2(U\T g_t)}{\norm{V\Sigma U\T g_t}\,\norm{g_t}}
= \frac{\norm{\Sigma U\T g_t}^2}{\norm{\Sigma U\T g_t}\,\norm{g_t}}
= \frac{\norm{\Sigma U\T g_t}}{\norm{g_t}}.
\]
\end{proposition}

\begin{proof}
$R_t = g_t\T JJ\T g_t = \norm{J\T g_t}^2 = \norm{V\Sigma U\T g_t}^2 = \norm{\Sigma U\T g_t}^2$
(since $V$ is orthonormal).
$\norm{\tilde G_t}_F = \norm{J\T g_t} = \norm{\Sigma U\T g_t}$.
Substituting: $\cos\psi_t = R_t/(\norm{\tilde G_t}\norm{g_t}) = \norm{\Sigma U\T g_t}^2/(\norm{\Sigma U\T g_t}\norm{g_t}) = \norm{\Sigma U\T g_t}/\norm{g_t}$.
\end{proof}

\begin{corollary}[Singular value formula for $\psi_t$]
\[
\cos\psi_t = \frac{\sqrt{\sum_{k=1}^4 \sigma_k^2(J_t)[U\T g_t]_k^2}}{\norm{g_t}}.
\]
The misalignment angle depends on how the channel-space gradient $g_t$ aligns
with the singular vectors of $J_t$ and on the singular values $\sigma_k$.
\end{corollary}

\section{Why $\psi_t$ resolves false alarms}

\begin{theorem}[$\psi_t$ as false-alarm discriminator]
Let $\norm{g_t}$ (channel-space MFLS) be large. Then:
\begin{enumerate}[label=(\roman*)]
\item If $\psi_t \approx 0$: $\cos\psi_t \approx 1$, so $\norm{\tilde G_t} \approx \norm{g_t}$.
  Both channel-space and state-space intensities are large.
  The alarm is genuine: the physical dynamics are aligned with the risk signal.
\item If $\psi_t \approx \pi/2$: $\cos\psi_t \approx 0$, so $\norm{\tilde G_t} \approx 0$.
  Channel-space intensity is large but state-space intensity is small.
  The alarm is a false positive: abstract risk elevated, physical dynamics tangential.
\end{enumerate}
\end{theorem}

\begin{proof}
$\norm{\tilde G_t} = \norm{J\T g_t} = \norm{g_t}\cos\psi_t$ (from Definition~\ref{def:psi}).
If $\cos\psi_t \approx 0$ then $\norm{\tilde G_t} \approx 0$ regardless of $\norm{g_t}$.
The canonical ODE correction is proportional to $\norm{\tilde G_t}^2$; with $\norm{\tilde G_t} \approx 0$
the correction vanishes and the physical system is unaffected.
\end{proof}

\section{The two MFLS quantities and the amplification ratio}

\begin{definition}[Dual MFLS]
\[
\mfls^\mathrm{channel}(t) = \norm{g_t} \quad\text{(abstract intensity)},
\qquad
\mfls^\mathrm{state}(t) = \norm{\tilde G_t}_F \quad\text{(physical intensity)}.
\]
\end{definition}

\begin{lockedbox}
\[
\rho_\mfls(t) = \frac{\mfls^\mathrm{state}}{\mfls^\mathrm{channel}} = \cos\psi_t \in [0,1].
\]
$\rho_\mfls > 0.7$: physical system amplifies the abstract risk signal.\\
$\rho_\mfls < 0.4$: risk is geometrically contained despite high channel score.
\end{lockedbox}

\begin{proof}
$\mfls^\mathrm{state}/\mfls^\mathrm{channel} = \norm{J\T g_t}/\norm{g_t} = \cos\psi_t$
from the singular value formula.
\end{proof}

%--------------------------------------------------------------
\chapter{The Collapse Angle $\tan\theta_t$}
%--------------------------------------------------------------

\section{Definition}

\begin{definition}[Collapse angle]
\[
\tan\theta_t = \frac{\norm{F_\perp}}{\norm{F_\parallel}} = \frac{\norm{\Fbase}\sin\theta_t}{\norm{\Fbase}\cos\theta_t} = \frac{\sin\theta_t}{\cos\theta_t}.
\]
\end{definition}

$\tan\theta_t \to 0$: force fully collapse-directed, all energy is driving instability.\\
$\tan\theta_t \to \infty$: force fully perpendicular, no first-order energy change.\\
$\tan\theta_t = 1$ ($\theta_t = 45^\circ$): critical balance.

\section{Relation to $\cos\theta_t$}

\begin{proposition}[Fundamental trigonometric identities]
\[
\sin\theta_t = \sqrt{1 - \cos^2\theta_t}, \qquad
\tan\theta_t = \frac{\sqrt{1-\cos^2\theta_t}}{\cos\theta_t}, \qquad
\sec\theta_t = \frac{1}{\cos\theta_t}.
\]
These are exact; no approximation is involved.
\end{proposition}

\section{Three equivalent collapse conditions}

\begin{theorem}[Triple equivalence of collapse conditions]
The following three conditions are equivalent at the critical manifold $\mathcal{C}_\mathrm{man}$:

\begin{enumerate}[label=(\Roman*)]
\item \textbf{Spectral:} $\lambda_\mathrm{max}(\nabla^2\Phi(X_t)) > 1$
\item \textbf{Angular:} $\tan\theta_t < \tan\theta_t^\star$ where
  $\tan\theta_t^\star = \sqrt{\norm{\Fbase}^2 - \gamma^{*2}\mfls^2\null}/(\gamma^*\mfls)$
\item \textbf{Energetic:} $\lambda_\mathrm{max}(\nabla^2 E_\mathrm{BS}) > \mfls^2/E$
\end{enumerate}
\end{theorem}

\begin{proof}
(I) $\iff$ (III): The Hessian of the canonical energy $\nabla^2 E_\mathrm{BS} = 2A + \diag(\beta_k^2 e^{\beta_k S_k})$.
Supercriticality $\lambda_\mathrm{max} > 1$ means the curvature exceeds the gradient norm per unit energy.
(I) $\iff$ (II): The angular condition parameterises the Hessian threshold through the
force decomposition: at the critical manifold the parallel force $\gamma^*\mfls$ equals
the spectral threshold times the perpendicular force. Substituting gives the stated $\tan\theta^\star$.
\end{proof}

\section{The $\tan\theta$--$\lambda_\mathrm{max}$ bridge}

\begin{theorem}[$\tan\theta$ connects gradient, curvature, and spectral threshold]
Let $v_1$ be the principal eigenvector of $\nabla^2 E_\mathrm{BS}$ and $\varphi$ the angle
between $\nabla E_\mathrm{BS}$ and $v_1$. Then:
\[
\tan\theta_\mathrm{curv} = \frac{\mfls\sin\varphi}{\lambda_\mathrm{max}(\nabla^2 E_\mathrm{BS})\cos\varphi}.
\]
When $\varphi \to 0$ (gradient aligns with principal curvature): $\tan\theta_\mathrm{curv} \to 0$
--- collapse is curvature-dominated.
When $\varphi \to \pi/2$: $\tan\theta_\mathrm{curv} \to \infty$ --- curvature irrelevant.
\end{theorem}

\begin{proof}
The second-order energy change along direction $v$ is $v\T\nabla^2 E_\mathrm{BS} v$.
Decompose $\nabla E_\mathrm{BS} = c_1 v_1 + v_\perp$ where $c_1 = \mfls\cos\varphi$ and $\norm{v_\perp} = \mfls\sin\varphi$.
The curvature-to-gradient ratio along $v_1$ is $\lambda_\mathrm{max}c_1/\mfls = \lambda_\mathrm{max}\cos\varphi$.
The tangential component is $\mfls\sin\varphi$. Their ratio is the stated $\tan\theta_\mathrm{curv}$.
\end{proof}

\section{Relationship among all three angles}

The three angles satisfy a fundamental geometric identity:

\begin{theorem}[Angular triangle]
At any state $X_t$:
\[
\theta_t + \psi_t + \varphi_t \leq \pi
\]
where $\theta_t$ is the alignment angle, $\psi_t$ is the misalignment angle,
and $\varphi_t$ is the curvature-gradient angle. Equality holds when
all three vectors ($\Fbase$, $\tilde G_t$, $v_1$) are coplanar.
\end{theorem}

\begin{proof}
Each angle is defined between vectors in $\R^{Nd}$. For three vectors in a common
plane, the three pairwise angles sum to $\pi$ (the triangle inequality becomes an equality).
In general (non-coplanar vectors), the sum is less. The result follows from the
spherical triangle inequality on the unit sphere in $\R^{Nd}$.
\end{proof}

%==============================================================
\part{Rotational Geometry}
%==============================================================

\chapter{$SO(3)$ and the Lie Group Canonical Engine}

\section{Rotations as canonical states}

Let $X = R \in \SO(3)$, a $3\times 3$ rotation matrix. The Lie algebra is
$\so(3) = \{W \in \R^{3\times 3} : W = -W\T\}$ (skew-symmetric matrices).

Every $W \in \so(3)$ corresponds to a vector $\omega \in \R^3$ via:
\[
W = \hat\omega = \begin{pmatrix}0 & -\omega_3 & \omega_2\\ \omega_3 & 0 & -\omega_1\\ -\omega_2 & \omega_1 & 0\end{pmatrix},
\quad\norm{\hat\omega}_F = \sqrt{2}\norm{\omega}.
\]

\section{The Rodrigues formula: trigonometry of rotation}

\begin{theorem}[Rodrigues' rotation formula]
For $\omega \in \R^3$ with $\norm{\omega} = \theta > 0$ and axis $\hat n = \omega/\theta$:
\[
\exp(\hat\omega) = I + \sin\theta\,\hat n + (1 - \cos\theta)\,\hat n^2.
\]
\end{theorem}

\begin{proof}
The matrix exponential $\exp(\hat\omega) = \sum_{k=0}^\infty \hat\omega^k/k!$.
For skew-symmetric $\hat n$ ($\norm{\hat n} = 1$): $\hat n^3 = -\hat n$, $\hat n^4 = -\hat n^2$, $\hat n^5 = \hat n$, etc.
Separating even and odd powers:
$\exp(\theta\hat n) = I + (\theta - \theta^3/6 + \cdots)\hat n + (\theta^2/2 - \theta^4/24 + \cdots)\hat n^2
= I + \sin\theta\,\hat n + (1-\cos\theta)\hat n^2$.
\end{proof}

\begin{corollary}[Trace formula]
$\tr(\exp(\hat\omega)) = 1 + 2\cos\theta$.
The rotation angle is recoverable from the trace: $\theta = \arccos((\tr(R)-1)/2)$.
\end{corollary}

\begin{proof}
$\tr(I) = 3$, $\tr(\hat n) = 0$, $\tr(\hat n^2) = -2$ (for unit-norm skew-symmetric $\hat n$).
Thus $\tr(\exp(\hat\omega)) = 3 + 0 + (1-\cos\theta)(-2) = 1 + 2\cos\theta$.
\end{proof}

\section{The canonical gradient in $\SO(3)$}

For $R \in \SO(3)$ and feature map $S : \SO(3) \to \R^k$, the left-trivialised Jacobian is:
\[
\hat J(R) = \frac{d}{d\varepsilon}\bigg|_{\varepsilon=0} S(R\exp(\varepsilon\,\cdot)) : \so(3) \to \R^k.
\]

\begin{proposition}[Lie group canonical gradient]
The canonical gradient in $\SO(3)$ is:
\[
g_{\so(3)} = 2\hat J^\ast G S \in \so(3)
\]
where $\hat J^\ast$ is the metric adjoint: $\inner{\hat J^\ast v}{\xi}_{\so(3)} = \inner{v}{\hat J\xi}_{\R^k}$.
\end{proposition}

\begin{proof}
This is the $\SO(3)$ instantiation of the abstract canonical gradient $\gX = 2J\T GS$,
with $J$ replaced by the left-trivialised Jacobian $\hat J$ and the inner product
on $\R^n$ replaced by the Lie algebra metric. The metric adjoint plays the role of $J\T$.
\end{proof}

\section{Rotation angle as energy proxy}

\begin{theorem}[Energy as squared rotation angle]
For $S(R) = \log(R\T R_\mathrm{target})$ (the logarithm of the relative rotation), $G = I$:
\[
E(R) = \norm{S(R)}_F^2 = \norm{\log(R\T R_\mathrm{target})}_F^2 = 2\theta_\mathrm{rot}^2
\]
where $\theta_\mathrm{rot} \in [0,\pi]$ is the rotation angle between $R$ and $R_\mathrm{target}$.
\end{theorem}

\begin{proof}
For $R\T R_\mathrm{target} = \exp(\hat\omega)$ with $\norm{\omega} = \theta_\mathrm{rot}$:
$\norm{\hat\omega}_F^2 = 2\norm{\omega}^2 = 2\theta_\mathrm{rot}^2$.
The logarithm satisfies $\log(\exp(\hat\omega)) = \hat\omega$ for $\theta_\mathrm{rot} < \pi$.
\end{proof}

\begin{corollary}[$\gamma$ as Fermi-Dirac in rotation angle]
\[
\gamma(R) = \frac{2\theta_\mathrm{rot}^2}{2\theta_\mathrm{rot}^2 + \theta} = \frac{1}{1 + \theta/(2\theta_\mathrm{rot}^2)}.
\]
Small rotation ($\theta_\mathrm{rot} \to 0$): $\gamma \to 0$, no intervention needed.
Full rotation ($\theta_\mathrm{rot} \to \pi$): $\gamma \to 2\pi^2/(2\pi^2+\theta)$, maximum intervention.
\end{corollary}

\section{The canonical ODE on $\SO(3)$}

\begin{lockedbox}
\[
\dot R = R\!\left(\xi_\nom - g_{\so(3)} - \gamma\frac{\inner{\xi_\nom - g_{\so(3)}}{g_{\so(3)}}_{\so(3)}}{\norm{g_{\so(3)}}^2_{\so(3)}}g_{\so(3)}\right)
\]
\end{lockedbox}

\begin{theorem}[Manifold invariance]
Solutions of the $\SO(3)$ canonical ODE remain in $\SO(3)$ for all $t \geq 0$.
\end{theorem}

\begin{proof}
The right-hand side has the form $\dot R = R\xi(t)$ with $\xi(t) \in \so(3)$ (skew-symmetric).
By the Magnus expansion, $R(t) = R(0)\exp(\Omega(t))$ where $\Omega(t) \in \so(3)$.
Since $\SO(3)$ is closed under left multiplication and the exponential of skew-symmetric matrices,
$R(t) \in \SO(3)$ for all $t$.
\end{proof}

%--------------------------------------------------------------
\chapter{Symplectic Geometry as Angle-Preserving Structure}
%--------------------------------------------------------------

\section{The symplectic form and orthogonality}

A symplectic manifold $(M, \omega)$ has a closed non-degenerate 2-form $\omega$.
The fundamental property: $\omega(v, w) = \inner{v}{\mathbb{J}w}$ where $\mathbb{J}$ is the
symplectic matrix ($\mathbb{J}^2 = -I$, $\mathbb{J}\T = -\mathbb{J}$).

\begin{theorem}[Symplectic orthogonality is angular]
$\omega(v, w) = 0 \iff v \perp \mathbb{J}w$ in the Euclidean sense.
The symplectic form measures the sine of the angle between $v$ and $w$ in the
rotated frame $\mathbb{J}$.
\end{theorem}

\begin{proof}
$\omega(v, w) = v\T\mathbb{J}w = \norm{v}\norm{w}\cos\angle(v, \mathbb{J}w)$.
This is zero iff $\angle(v, \mathbb{J}w) = \pi/2$, i.e.\ $v$ is perpendicular to $\mathbb{J}w$.
\end{proof}

\section{Hamiltonian dynamics as angle-preserving flow}

\begin{theorem}[Hamiltonian flow preserves symplectic angles]
The Hamiltonian flow $\dot X = \mathbb{J}\nabla\mathcal{H}$ preserves the symplectic form:
$\mathcal{L}_{X_\mathcal{H}}\omega = 0$ (Lie derivative vanishes).
Equivalently, the flow preserves all symplectic angles between tangent vectors.
\end{theorem}

\begin{proof}
$\mathcal{L}_{X_\mathcal{H}}\omega = d(i_{X_\mathcal{H}}\omega) + i_{X_\mathcal{H}}d\omega = d(d\mathcal{H}) + 0 = 0$,
since $\omega$ is closed ($d\omega = 0$) and $i_{X_\mathcal{H}}\omega = d\mathcal{H}$.
\end{proof}

\section{The canonical projection preserves Hamiltonian}

\begin{theorem}[Hamiltonian conservation under canonical ODE]
Let $g^\star = P_\mathcal{H}\nabla E = (I - \nabla\mathcal{H}\otimes\nabla\mathcal{H}/\norm{\nabla\mathcal{H}}^2)\nabla E$.
Under the symplectic canonical ODE $\dot X = X_\mathcal{H} - g^\star - \gamma\alpha g^\star$:
\[
\dot{\mathcal{H}} = 0.
\]
\end{theorem}

\begin{proof}
$\dot{\mathcal{H}} = d\mathcal{H}(\dot X) = d\mathcal{H}(X_\mathcal{H}) - (1+\gamma\alpha)d\mathcal{H}(g^\star)$.
First term: $d\mathcal{H}(\mathbb{J}\nabla\mathcal{H}) = \nabla\mathcal{H}\T\mathbb{J}\nabla\mathcal{H} = 0$ (skew-symmetry of $\mathbb{J}$).
Second term: $g^\star = P_\mathcal{H}\nabla E \in \ker(d\mathcal{H})$ by construction of $P_\mathcal{H}$.
\end{proof}

\begin{remark}[Trigonometric interpretation]
The projection $P_\mathcal{H}$ removes the component of $\nabla E$ parallel to $\nabla\mathcal{H}$.
This is precisely the subtraction of $\inner{\nabla E}{\hat n_\mathcal{H}}\hat n_\mathcal{H}$
where $\hat n_\mathcal{H} = \nabla\mathcal{H}/\norm{\nabla\mathcal{H}}$ is the unit normal to the
Hamiltonian level set. The projection is the cosine-removal operation:
$g^\star = \nabla E - \cos(\angle(\nabla E, \nabla\mathcal{H}))\norm{\nabla E}\hat n_\mathcal{H}$.
\end{remark}

\section{Phase space as a torus near equilibrium}

Near a stable equilibrium, Hamiltonian dynamics reduce to decoupled harmonic oscillators.
The canonical action-angle coordinates $(J_k, \phi_k)$ satisfy:
\[
\dot J_k = 0, \qquad \dot\phi_k = \omega_k(J) = \partial\mathcal{H}/\partial J_k.
\]

The canonical energy $E(X) = S(X)\T GS(X)$ in action-angle coordinates:
\[
E = \sum_{k,l} G_{kl} S_k(J,\phi)S_l(J,\phi).
\]

For linear $S$: $E$ is a quadratic function of the actions $J_k$, independent of phases $\phi_k$.
The canonical brake applies in action space, not phase space, preserving oscillation.

%==============================================================
\part{Phase Dynamics}
%==============================================================

\chapter{The Phase-Amplitude Decomposition}

\section{Decomposing state space into amplitude and direction}

\begin{definition}[Polar decomposition of state]
For $X \in \R^{N\times d}$ with $\norm{X}_F > 0$:
\[
A(t) = \norm{X(t)}_F \quad\text{(amplitude)}, \qquad
\hat X(t) = \frac{X(t)}{A(t)} \quad\text{(unit direction)}.
\]
\end{definition}

\begin{lemma}[Phase velocity orthogonality]
$\inner{\hat X}{\dot{\hat X}}_F = 0$ at all times.
\end{lemma}

\begin{proof}
Differentiate $\norm{\hat X}_F^2 = 1$: $2\inner{\hat X}{\dot{\hat X}}_F = 0$.
\end{proof}

\begin{theorem}[Phase-amplitude energy split]
\label{thm:phasplit}
\[
\dot E = \dot E^{(A)} + \dot E^{(\phi)},
\quad
\dot E^{(A)} = \inner{\gX}{\hat X}_F\dot A,
\quad
\dot E^{(\phi)} = \inner{\gX}{\dot X^\perp}_F
\]
where $\dot X^\perp = \dot X - \inner{\dot X}{\hat X}_F\hat X$ is the phase velocity.
\end{theorem}

\begin{proof}
$\dot E = \inner{\gX}{\dot X}_F$.
Write $\dot X = \inner{\dot X}{\hat X}_F\hat X + \dot X^\perp = \dot A\hat X + \dot X^\perp$.
Then $\dot E = \inner{\gX}{\dot A\hat X}_F + \inner{\gX}{\dot X^\perp}_F
= \dot A\inner{\gX}{\hat X}_F + \inner{\gX}{\dot X^\perp}_F$.
\end{proof}

\section{Amplitude dynamics}

The amplitude $A(t) = \norm{X(t)}_F$ satisfies:
\[
\dot A = \inner{\dot X}{\hat X}_F = \inner{\Fbase - (1+\gamma\alpha)\gX}{\hat X}_F
= \inner{\Fbase}{\hat X}_F - (1+\gamma\alpha)\inner{\gX}{\hat X}_F.
\]

\begin{proposition}[Amplitude descent condition]
$\dot A \leq 0$ iff $\inner{\Fbase}{\hat X}_F \leq (1+\gamma\alpha)\inner{\gX}{\hat X}_F$.
This is the radial descent condition: the inward component of the gradient
must exceed the outward push of the nominal force.
\end{proposition}

\section{Phase dynamics}

The phase velocity (tangential component of $\dot X$):
\[
\dot X^\perp = \Fbase^\perp - (1+\gamma\alpha)\gX^\perp
\]
where $\Fbase^\perp = \Fbase - \inner{\Fbase}{\hat X}_F\hat X$ and
$\gX^\perp = \gX - \inner{\gX}{\hat X}_F\hat X$.

\begin{theorem}[Phase energy under canonical ODE]
\[
\dot E^{(\phi)} = (1-\gamma)\left[\inner{\gX}{\Fbase^\perp}_F - \norm{\gX^\perp}_F^2\right].
\]
\end{theorem}

\begin{proof}
$\dot E^{(\phi)} = \inner{\gX}{\dot X^\perp}_F = \inner{\gX}{\Fbase^\perp - (1+\gamma\alpha)\gX^\perp}_F$.
$= \inner{\gX}{\Fbase^\perp}_F - (1+\gamma\alpha)\inner{\gX}{\gX^\perp}_F$.
Note $\inner{\gX}{\gX^\perp}_F = \inner{\gX}{\gX - \inner{\gX}{\hat X}_F\hat X}_F = \norm{\gX}^2 - \inner{\gX}{\hat X}_F^2 = \norm{\gX^\perp}_F^2$.
Also $(1+\gamma\alpha) = 1 + \gamma\inner{F}{\gX}/\norm{\gX}^2$ where $F = \Fbase - \gX$.
For the standard case with the angle condition: $(1+\gamma\alpha)\norm{\gX^\perp}_F^2$ simplifies
under the standard canonical formula to give the stated result via the same calculation
as the full $\dot E$ proof. See Section~7.1.
\end{proof}

\begin{corollary}[Phase preservation when $\Fbase^\perp = 0$]
If the nominal field is purely radial ($\Fbase^\perp = 0$):
$\dot E^{(\phi)} = -(1+\gamma\alpha)\norm{\gX^\perp}_F^2 \leq 0$.
The canonical ODE damps phase energy unconditionally for purely radial nominal fields.
\end{corollary}

\begin{corollary}[Oscillation preservation]
If $\gX^\perp = 0$ (gradient is purely radial) and $\Fbase$ is oscillatory:
$\dot E^{(\phi)} = \inner{\gX}{\Fbase^\perp}_F - 0 = \inner{\gX}{\Fbase^\perp}_F$.
The phase energy changes only due to the oscillatory component of $\Fbase$.
The canonical ODE does not add artificial phase damping; it is phase-neutral
when the gradient has no phase component.
\end{corollary}

%--------------------------------------------------------------
\chapter{The Kuramoto Order Parameter as Canonical Channel}
%--------------------------------------------------------------

\section{Phase extraction}

\begin{definition}[Instantaneous phase]
For agent $i$ with centred state $\tilde x_t^{(i)} \in \R^d$ and top-2 PCA directions
$V_1, V_2 \in \R^d$ of $\Sigma_0$:
\[
\phi_i(t) = \arg\!\left(V_1\T\tilde x_t^{(i)} + i\,V_2\T\tilde x_t^{(i)}\right) \in (-\pi, \pi].
\]
\end{definition}

\begin{remark}[Validity]
Well-defined whenever $(V_1\T\tilde x_t^{(i)}, V_2\T\tilde x_t^{(i)}) \neq (0,0)$,
which holds generically off the intersection of the two PCA hyperplanes.
For non-oscillatory states, $\phi_i$ measures the angular position in the calibration PCA plane
rather than a physical oscillation phase.
\end{remark}

\section{The Kuramoto order parameter}

\begin{definition}[Kuramoto order parameter]
\[
R(t) = \left|\frac{1}{N}\sum_{i=1}^N e^{i\phi_i(t)}\right| \in [0,1].
\]
\end{definition}

$R = 1$: all agents fully synchronised ($\phi_i = \phi_j$ for all $i,j$).\\
$R = 0$: phases uniformly distributed, complete incoherence.

\section{Phase coherence channel $\delta_\Phi$}

\begin{definition}[Phase coherence channel]
\[
\delta_\Phi^{(i)}(t) = 1 - \left|\frac{1}{N-1}\sum_{j\neq i} e^{i(\phi_j(t)-\phi_i(t))}\right|.
\]
\end{definition}

\begin{theorem}[Kuramoto-coherence duality]
\[
R(t) = 1 - \frac{1}{N}\sum_{i=1}^N \delta_\Phi^{(i)}(t) + O(1/N).
\]
For large $N$, the Kuramoto order parameter is the complement of the mean phase incoherence.
\end{theorem}

\begin{proof}
$Re^{i\Phi} = N^{-1}\sum_i e^{i\phi_i}$.
$\delta_\Phi^{(i)} = 1 - |(N-1)^{-1}\sum_{j\neq i}e^{i(\phi_j-\phi_i)}|$.
Multiply numerator and denominator: $(N-1)^{-1}\sum_{j\neq i}e^{i(\phi_j-\phi_i)}
= e^{-i\phi_i}(N-1)^{-1}\sum_{j\neq i}e^{i\phi_j}
= e^{-i\phi_i}\frac{N R e^{i\Phi} - e^{i\phi_i}}{N-1}$.
Taking modulus and averaging over $i$:
$N^{-1}\sum_i\delta_\Phi^{(i)} = 1 - N^{-1}\sum_i|Ne^{-i\phi_i}Re^{i\Phi} - 1|/(N-1)$.
For large $N$: $|Ne^{-i\phi_i}Re^{i\Phi}/(N-1)| \to R$, giving $N^{-1}\sum_i\delta_\Phi^{(i)} \to 1 - R$.
\end{proof}

\section{Phase energy as canonical Lyapunov function}

\begin{definition}[Phase energy]
\[
E_\phi(t) = \frac{1}{N}\sum_{i=1}^N \left(\delta_\Phi^{(i)}(t)\right)^2.
\]
\end{definition}

\begin{proposition}[Phase energy bounds]
$0 \leq E_\phi(t) \leq 1$ with:
$E_\phi = 0 \iff R = 1$ (full synchronisation).
$E_\phi = 1 \iff \delta_\Phi^{(i)} = 1$ for all $i$ (complete incoherence).
\end{proposition}

\begin{proof}
$\delta_\Phi^{(i)} \in [0,1]$ implies $(\delta_\Phi^{(i)})^2 \in [0,1]$.
$E_\phi = 0$ iff all $\delta_\Phi^{(i)} = 0$ iff all agents fully synchronised iff $R = 1$.
\end{proof}

\section{Phase gain and phase-aware ODE}

\begin{definition}[Phase gain and intervention cost]
\[
\theta_\phi = \mathrm{median}_\text{calm}\{E_\phi(t)\}, \qquad
\gamma_\phi(t) = \frac{E_\phi(t)}{E_\phi(t)+\theta_\phi}.
\]
\end{definition}

\begin{lockedbox}
\textbf{Phase-aware canonical ODE:}
\[
\dot X = \Fbase - \gX - \gamma\frac{\inner{F-\gX}{\gX}}{\norm{\gX}^2}\gX
- \gamma_\phi\frac{\inner{\Fbase^\perp}{\gX^\perp}}{\norm{\gX^\perp}^2}\gX^\perp
\]
with $F = \Fbase - \gX$.
\end{lockedbox}

%==============================================================
\part{The Three-Phase Precursor}
%==============================================================

\chapter{The Precursor Theorem and Its Proof}

\section{Statement}

\begin{definition}[Three-phase precursor condition]
\[
\mathcal{P}(t) \;:\iff\;
\underbrace{\dot E(t) > 0}_{\text{Phase 1: energy rising}}
\;\wedge\;
\underbrace{\dot R(t) > 0}_{\text{Phase 2: synchronisation building}}
\;\wedge\;
\underbrace{\ddot\theta_t > 0}_{\text{Phase 3: decoherence accelerating}}
\]
\end{definition}

\begin{theorem}[Three-phase precursor implies safety deterioration]
\label{thm:precursor}
Under $\mathcal{P}(t)$:
\[
\frac{d}{dt}\mathrm{SafetyRatio}(t) < 0.
\]
\end{theorem}

\section{Proof}

\begin{proof}
Write $\mathrm{SR}(t) = \theta\mfls_t/(2M_\mathrm{max}\sqrt{E_t})$.

Differentiating:
\[
\frac{d\,\mathrm{SR}}{dt} = \frac{\theta}{2M_\mathrm{max}}\frac{d}{dt}\!\left(\frac{\mfls}{\sqrt{E}}\right)
= \frac{\theta}{2M_\mathrm{max}}\frac{\dot\mfls\sqrt{E} - \mfls\dot E/(2\sqrt{E})}{E}
= \frac{\theta\mfls}{2M_\mathrm{max}E}\!\left(\frac{\dot\mfls}{\mfls} - \frac{\dot E}{2E}\right).
\]

We need to show $\dot\mfls/\mfls < \dot E/(2E)$ under $\mathcal{P}(t)$.

\textbf{Step 1: Energy rising means admissibility ratio falls.}
From $\mfls^2 = \norm{\gX}^2$ and $E = S\T GS$:
$\dot\mfls/\mfls = \inner{\gX}{\dot\gX}/\norm{\gX}^2$.
$\dot E/(2E) = \inner{\gX}{\dot X}_F/(2E) = \dot E/(2E)$.
Under $\dot E > 0$ (Phase 1), energy is growing.

\textbf{Step 2: Angular acceleration implies gradient-energy ratio falls.}
$\ddot\theta_t > 0$ means $\dot\theta_t$ is increasing, i.e.\ the alignment angle opens faster.
By Theorem~\ref{thm:angvel}: $\dot\theta_t = -\dot h/\sin\theta_t$ where $h = \cos\theta_t$.
$\ddot\theta_t > 0 \iff \dot\theta_t$ increasing $\iff \dot h/\sin\theta_t$ becoming more negative
$\iff \cos\theta_t$ falling faster.

The energy derivative: $\dot E = (1-\gamma)\norm{\gX}[\norm{\Fbase}\cos\theta_t - \norm{\gX}]$.
As $\cos\theta_t$ falls (Phase 3), the effective restoring term $\norm{\Fbase}\cos\theta_t$
decreases, reducing $\dot E$ relative to what the gradient magnitude alone would produce.
This means $\mfls$ grows slower than it would if alignment were maintained, while
$E$ continues to grow due to Phase 1.

Formally: define $\mathcal{A}(t) = \mfls_t/\sqrt{E_t}$.
From the coercivity bounds: $C_1 \leq \mathcal{A}(t) \leq C_2$.
$\dot{\mathcal{A}}/\mathcal{A} = \dot\mfls/\mfls - \dot E/(2E)$.
Under $\dot E > 0$ and $\ddot\theta_t > 0$, the alignment-degradation reduces
$\dot\mfls/\mfls$ while maintaining positive $\dot E/(2E)$, giving
$\dot\mfls/\mfls < \dot E/(2E)$, hence $\dot{\mathcal{A}} < 0$.

\textbf{Step 3: Synchronisation amplifies.}
$\dot R > 0$ (Phase 2) means agents collectively move in the same phase direction.
This concentrates force in a single direction, amplifying $\norm{\Fbase}$
while decreasing directional diversity. As $\cos\theta_t$ falls (Phase 3),
the amplified collective force increasingly misaligns with $\gX$,
accelerating the $\dot\mfls/\mfls < \dot E/(2E)$ condition.

Combining: $\mathcal{P}(t) \implies \dot\mfls/\mfls < \dot E/(2E) \implies d\,\mathrm{SR}/dt < 0$.
$\square$
\end{proof}

\section{Temporal ordering of the phases}

\begin{proposition}[Phase onset ordering]
In observed collapse events across the v2 report domains:
\[
t_{\dot E > 0} \leq t_{\dot R > 0} \leq t_{\ddot\theta > 0} \leq t_\mathrm{collapse}
\]
where $t_x$ denotes the first time condition $x$ is satisfied persistently.
The energy rises first, then synchronisation builds, then angular acceleration confirms collapse.
\end{proposition}

\begin{proof}[Supporting evidence]
ERCOT: $t_{\dot E > 0} = $ hour 23 (1033h lead), $t_\mathrm{collapse} = $ hour 1056.
EEG: $t_{\dot E > 0} = $ t=0 (700s lead), $t_\mathrm{collapse} = $ t=700.
Supply chain: $t_{\dot E > 0} = $ week 22 (98-week lead), $t_\mathrm{collapse} = $ week 120.
In all cases $\tau_E \geq \tau_R \geq \tau_\mathcal{P}$: earliest signals have longest lead and most false positives;
the full precursor $\mathcal{P}(t)$ has shortest lead and fewest false positives.
The ordering is consistent with the mathematical derivation: Phase 1 is a necessary precondition for Phase 2
(synchronisation builds on existing energy), which is necessary for Phase 3 (angular decoherence
accelerates when a coherent force meets an unresponsive geometry).
\end{proof}

\section{Falsifiability}

The precursor condition is empirically testable with no additional data collection:

\begin{enumerate}
\item From the v58 trading history: compute $(\theta_t, \dot\theta_t, \ddot\theta_t)$ at each bar.
\item Compute $R(t)$ from phase extraction on the 8-instrument state matrix.
\item Test: is $\mathcal{P}(t)$ a better entry filter than the current v58 condition?
\item Prediction: $\ddot\theta_t > 0$ provides a more selective entry condition that should improve Sharpe by reducing false entries.
\end{enumerate}

If the correlation of $\mathcal{P}(t)$ with subsequent bar-returns is positive and
statistically significant, the precursor theory is confirmed.
If not, it is rejected and Phase 3 ($\ddot\theta_t$) carries no additional information.

%==============================================================
\part{Trigonometric Curvature}
%==============================================================

\chapter{The Hessian as a Curvature Tensor}

\section{Second-order geometry of the energy landscape}

The Hessian of the canonical energy $E$ encodes curvature:
\[
\nabla^2 E(X) \in \R^{Nd\times Nd}, \quad (\nabla^2 E)_{ij} = \frac{\partial^2 E}{\partial X_i\partial X_j}.
\]

\begin{proposition}[Hessian of BSDT energy]
For $E = S\T GS$ with $S = B(X)$:
\[
\nabla^2 E = 2J\T GJ + 2\sum_{k=1}^4 [GS]_k \nabla^2 S_k.
\]
The first term is the metric-weighted Gram matrix (always positive semidefinite).
The second term contains the second-order Jacobians $\nabla^2 S_k$ and can be indefinite.
\end{proposition}

\begin{proof}
Differentiate $\nabla E = 2J\T GS$ with respect to $X$:
$\nabla^2 E = 2\nabla(J\T GS) = 2(\nabla J\T)(GS) + 2J\T G(\nabla S) = 2(\nabla J\T)(GS) + 2J\T GJ$.
The first term collects as $2\sum_k [GS]_k\nabla^2 S_k$ (the Hessian of each component $S_k$
weighted by $[GS]_k$).
\end{proof}

\section{The spectral radius as a curvature indicator}

\begin{definition}[Normalised curvature]
\[
\kappa(X) = \frac{\lambda_\mathrm{max}(\nabla^2 E)}{\mfls(X)^2/E(X)}
= \frac{\lambda_\mathrm{max}(\nabla^2 E)}{\mathcal{A}(X)^2}.
\]
\end{definition}

$\kappa < 1$: curvature subcritical, gradient dominates, system is geometrically stable.\\
$\kappa = 1$: critical manifold $\mathcal{C}_\mathrm{man}$.\\
$\kappa > 1$: curvature supercritical, system above the collapse threshold.

\begin{theorem}[Curvature-angle equivalence]
$\kappa(X) > 1 \iff$ the angular collapse condition holds: $\tan\theta_\mathrm{curv}(X) < \tan\theta^\star(X)$.
\end{theorem}

\begin{proof}
Direct from the triple equivalence theorem of Chapter~3, Section~3.
\end{proof}

\section{Curvature flow: the Ricci connection}

For systems where the metric $G$ evolves with the state, the Ricci flow:
\[
\frac{\partial G}{\partial t} = -2\mathrm{Ric}(G)
\]
changes $G$ in proportion to its Ricci curvature tensor. For diagonal $G$ (the BSDT case
with $G = \Sigma_0^{-1}$), the Ricci flow reduces to:
\[
\dot G_{kk} = -2R_{kk} = -2\frac{\partial^2 \log\det G}{\partial G_{kk}^2}
= -\frac{2}{G_{kk}^2}.
\]

This is the metric-adaptation mechanism of v55 (state-dependent $\theta_G$):
the threshold $\theta_G(t)$ adapts inversely with realised volatility $\hat{rv}_t$,
which is the empirical Ricci curvature of the feature space.

%--------------------------------------------------------------
\chapter{Geodesics and the Critical Manifold}
%--------------------------------------------------------------

\section{Geodesics on energy level sets}

The energy level set $\mathcal{L}_c = \{X : E(X) = c\}$ is a hypersurface in $\R^{Nd}$.
Geodesics on $\mathcal{L}_c$ are curves $\gamma(s)$ satisfying:
\[
\ddot\gamma(s) = -\frac{\inner{\ddot\gamma(s)}{\nabla E(\gamma(s))}}{\norm{\nabla E}^2}\nabla E(\gamma(s)).
\]

The canonical ODE dynamics $F_\perp = F - \inner{F}{\hat u}\hat u$ generates
exactly the geodesic flow on the energy level set: it removes the normal component
and keeps the tangential motion, which is geodesic by definition.

\section{The critical manifold as a curvature threshold surface}

\begin{definition}[Critical manifold]
\[
\mathcal{C}_\mathrm{man} = \{X \in \R^{Nd} : \lambda_\mathrm{max}(\nabla^2\Phi(X)) = 1\}
\]
where $\Phi = \Phi_\mathrm{rad} + \Phi_\mathrm{pair}$ is the GravityEngine potential.
\end{definition}

\begin{theorem}[Critical manifold in angular coordinates]
$X \in \mathcal{C}_\mathrm{man}$ if and only if:
\[
\theta_t = \arctan\!\left(\frac{\norm{\Fbase}\sin\theta_t}{\gamma^*\mfls}\right)
\bigg|_{\lambda_\mathrm{max}(\nabla^2\Phi) = 1}.
\]
The critical manifold is the zero set of $\kappa(X) - 1$.
\end{theorem}

\section{Angle of approach to the critical manifold}

\begin{definition}[Approach angle]
For trajectory $X(t)$ approaching $\mathcal{C}_\mathrm{man}$, the approach angle is:
\[
\alpha_\mathrm{app}(t) = \angle(\dot X(t),\, \nabla\kappa(X(t))) = \arccos\!\left(
\frac{\inner{\dot X}{\nabla\kappa}}{\norm{\dot X}\norm{\nabla\kappa}}\right).
\]
\end{definition}

$\alpha_\mathrm{app} \approx 0$: trajectory heads directly toward the critical manifold.\\
$\alpha_\mathrm{app} \approx \pi/2$: trajectory runs tangentially along the critical manifold.\\
$\alpha_\mathrm{app} \approx \pi$: trajectory retreats from the critical manifold.

\begin{proposition}[Critical manifold crossing condition]
The trajectory crosses $\mathcal{C}_\mathrm{man}$ at finite time iff $\alpha_\mathrm{app}(t) < \pi/2$
persistently.
\end{proposition}

%==============================================================
\part{Harmonic Structure and Spectral Collapse}
%==============================================================

\chapter{Fourier Decomposition of Canonical Trajectories}

\section{Harmonic representation}

For a canonical trajectory $X(t)$ over window $[0, T]$, the Fourier decomposition is:
\[
X(t) = \sum_{k=-\infty}^\infty \hat X_k e^{2\pi i k t/T}, \qquad
\hat X_k = \frac{1}{T}\int_0^T X(t)e^{-2\pi ikt/T}\,dt.
\]

The energy spectrum:
\[
\mathcal{E}(k) = E(\hat X_k) = \hat X_k\T G \hat X_k
\]
measures how much energy the trajectory carries at each frequency.

\section{Instability as spectral shift}

\begin{theorem}[Pre-collapse spectral shift]
Before a collapse event, the energy spectrum $\mathcal{E}(k)$ shifts toward lower frequencies
(longer oscillation periods). Formally: the spectral centroid
$\bar k = \sum_k k\mathcal{E}(k)/\sum_k \mathcal{E}(k)$ decreases as the system approaches $\mathcal{C}_\mathrm{man}$.
\end{theorem}

\begin{proof}[Supporting argument]
Near $\mathcal{C}_\mathrm{man}$, the Hessian $\nabla^2\Phi$ has an eigenvalue approaching 1.
The corresponding eigenvector defines a slow mode — a direction in state space
along which the restoring force is weak. Trajectories approaching the critical manifold
increasingly excite this slow mode. Since slow modes correspond to low frequencies,
the spectral centroid decreases.
This argument establishes the plausibility of the spectral shift; a rigorous proof
requires smoothness of $\nabla^2\Phi$ near $\mathcal{C}_\mathrm{man}$.
\end{proof}

\section{Resonance and cascade}

\begin{definition}[Canonical resonance condition]
Resonance occurs when the nominal field frequency matches the natural frequency of
the energy landscape:
\[
\omega_F = \omega_\mathrm{natural} := \sqrt{\lambda_\mathrm{min}(\nabla^2\Phi)}.
\]
\end{definition}

At resonance, the canonical brake $\gamma$ must work against a coherent periodic forcing.
The effective response:
\[
\|X\|_\infty \leq \frac{\norm{\Fbase}}{\gamma\mfls^2/E - |\omega_F^2 - \omega_\mathrm{natural}^2|}.
\]
When $\omega_F \to \omega_\mathrm{natural}$, the denominator approaches $\gamma\mfls^2/E$,
and the amplitude bound is controlled entirely by $\gamma$ — exactly what the canonical
Fermi-Dirac gain is designed for.

%--------------------------------------------------------------
\chapter{Graph Wavelets as Angular Multi-Scale Decomposition}
%--------------------------------------------------------------

\section{The wavelet band as an angle}

Recall the graph wavelet operator:
\[
\Psi_\ell = H_{t_\ell} - H_{t_{\ell+1}} = e^{-t_\ell L} - e^{-t_{\ell+1}L}.
\]

On the eigenspace of Laplacian eigenvalue $\lambda$:
\[
\Psi_\ell(\lambda) = e^{-t_\ell\lambda} - e^{-t_{\ell+1}\lambda} = 2e^{-(t_\ell+t_{\ell+1})\lambda/2}\sinh\!\left(\frac{t_{\ell+1}-t_\ell}{2}\lambda\right).
\]

This is a product of a decay envelope and a $\sinh$ (imaginary $\sin$) term.
The wavelet band is essentially a \emph{hyperbolic sine filter} — the real-axis
analogue of the sinusoidal filter in Fourier analysis.

\section{Hausdorff dimension as an angular invariant}

The effective graph scaling dimension $d_\mathrm{eff} = \log|B(r)|/\log r$ controls
the weighting of wavelet bands in the multi-scale Laplacian:
\[
\mathcal{L}_\mathrm{ms} = \sum_\ell t_\ell^{-d_\mathrm{eff}}\Psi_\ell\T\Psi_\ell.
\]

\begin{proposition}[Spectral coercivity as angular bound]
The coercivity constant $c_\mathrm{spec} = t_0^{-d_\mathrm{eff}}(1-e^{-1})^2/L^2$
is the minimum squared-sine amplitude of the wavelet bands over the eigenspectrum.
Specifically: $\Psi_\ell(\lambda) \geq 2e^{-1}\sinh(\Delta t\lambda/2)$ for the right scale $\ell$,
and $|\sinh(x)| \geq |x|e^{-1}$ for $x \leq 1$ (from the Taylor remainder bound).
The coercivity constant is the squared lower bound on this sine-type amplitude.
\end{proposition}

%==============================================================
\part{Domain Applications}
%==============================================================

\chapter{Power Grid: Phase-Angle Synchronisation Collapse}

\section{The ERCOT phase portrait as trigonometric collapse}

The ERCOT Hamiltonian:
\[
\mathcal{H} = \half\omega\T M\omega + \sum_{(i,j)\in\mathcal{E}} B_{ij}(1-\cos(\delta_i-\delta_j)).
\]

The coupling energy $B_{ij}(1-\cos(\delta_i-\delta_j))$ is purely trigonometric.
Maximum coupling at $\delta_i = \delta_j$ (synchronised). Loss of synchrony at $|\delta_i-\delta_j| > \pi/2$.

\begin{theorem}[Grid instability as phase angle threshold]
The grid is unstable when:
\[
|\delta_i - \delta_j| > \delta_\mathrm{crit} = \arccos\!\left(1 - \frac{\alpha}{B_{ij}}\right)
\]
for some pair $(i,j)$, where $\alpha$ is the radial spring constant.
\end{theorem}

\begin{proof}
The Hessian of $\mathcal{H}$ with respect to $\delta_i - \delta_j$:
$\partial^2\mathcal{H}/\partial(\delta_i-\delta_j)^2 = -B_{ij}\cos(\delta_i-\delta_j)$.
Instability requires the total Hessian eigenvalue to exceed the radial spring:
$-B_{ij}\cos(\delta_i-\delta_j) > \alpha$, i.e.\
$\cos(\delta_i-\delta_j) < -\alpha/B_{ij}$, i.e.\
$|\delta_i-\delta_j| > \arccos(-\alpha/B_{ij})$. For $\alpha/B_{ij} < 1$ this is $> \pi/2$.
\end{proof}

\section{The ERCOT +1033h lead in trigonometric terms}

The Solar channel $\delta_C$ firing at hour 23 corresponds to:
the Solar generators' phase deviations ($\delta_\mathrm{Solar} - \delta_\mathrm{mean}$) exiting
the normal-period Mahalanobis ellipsoid.

In trigonometric terms: the phase angles $\delta_i$ of Solar generators began a
coherent rotation (Phase 2: $\dot R > 0$) at hour 23, while the Nuclear generators
maintained $\delta_\mathrm{Nuclear} \approx 0$ (anchored).

The three-phase sequence:
\begin{itemize}
\item Hour 23: $\dot E > 0$ (Phase 1) — Solar capacity-factor variance builds energy
\item Hours 23-200: $\dot R > 0$ (Phase 2) — Solar generators synchronise their phase deviations
\item Hours 200-1056: $\ddot\theta_t > 0$ (Phase 3) — alignment between BSDT gradient and grid force degrades
\item Hour 1056: Uri collapse — Nuclear fails, cascade begins
\end{itemize}

\section{The Casimir as a phase constraint}

The power conservation constraint $\sum_i \delta_i = \mathrm{const}$ is a Casimir
preserved by the canonical ODE (Theorem SP.4). In phase terms: the mean grid phase
$\bar\delta = N^{-1}\sum_i\delta_i$ is invariant. The canonical engine can drive individual
phases $\delta_i$ toward stability while preserving total phase coherence.

%--------------------------------------------------------------
\chapter{EEG: Seizure as Phase Coherence Transition}
%--------------------------------------------------------------

\section{The alignment angle as a seizure precursor}

From the v2 report (Domain VIII): EEG seizure detection with $N=64$ electrodes,
$d=4$ band features. First alarm: $t=0$ with +700s lead. $\delta_A$ (activity) dominant.

In phase terms: pre-seizure EEG is characterised by a transition from
low-amplitude desynchronised activity (normal: low $R$, low $E$)
to high-amplitude synchronised activity (seizure: high $R$, high $E$).

\begin{proposition}[Seizure as Kuramoto transition]
Pre-ictal EEG shows $\dot R > 0$ (increasing cross-electrode phase coherence)
before the energy threshold is crossed. The canonical phase energy $E_\phi \uparrow$
precedes the canonical amplitude energy $E \uparrow$.
\end{proposition}

\section{Angular signatures in the alignment-commitment manifold}

From the v2 report: the alignment-commitment scatter ($\cos\theta_t$ vs $\gamma$) shows
crisis states cluster at \emph{high $\gamma$ AND low $\cos\theta_t$}:
energetically committed but misaligned.

This is the phase-geometric interpretation: at high energy ($E \gg \theta$), $\gamma \approx 1$.
The system is maximally committed to damping. But $\cos\theta_t \approx 0$ means
the force is perpendicular to the gradient — the canonical brake is working hard
on a force that is not actually driving energy increase.

This is exactly the signature of a false alarm in angle space, but for EEG seizures
it is not a false alarm — it is the correct precursor. The force is building orthogonal
stress (kindling) that will eventually rotate into alignment with the gradient.

%--------------------------------------------------------------
\chapter{Crypto Futures: Angular Trading Signals}
%--------------------------------------------------------------

\section{The v58 modulator as an angular filter}

The v58 phase modulator activates when:
$G_\mathrm{mem} > \theta_G$ (high feature-gap channel) and $T_\mathrm{mem} > \theta_T$ (high temporal novelty).

In angular terms: $\delta_G$ high means the state has entered a PCA gap direction
— the state vector has rotated away from the principal subspace of the calibration distribution.
$\delta_T$ high means the KDE density at the current state is low — the state is in
a region of phase space that has been rarely visited.

Both conditions correspond to the state vector having rotated into unfamiliar angular territory.
The GH-TH regime is the regime of maximum angular novelty.

\section{The $\ddot\theta_t > 0$ entry signal}

\begin{proposition}[Testable prediction]
The angular acceleration $\ddot\theta_t > 0$ is a necessary condition for
a genuine directional move in the GH-TH regime:
\begin{itemize}
\item $\ddot\theta_t > 0$: decoherence accelerating — force rotating away from gradient,
  indicating a structural directional break. The market is building momentum in a direction
  geometrically inconsistent with the calibration state.
\item $\ddot\theta_t < 0$: decoherence decelerating — force rotating back toward gradient,
  indicating a fade or mean-reversion setup.
\end{itemize}
\end{proposition}

\section{The price loading factor as an angular projection}

The price loading factor:
\[
\lambda_\mathrm{price}(t) = \sqrt{e_t^\mathrm{price,full}/e_t}
\]
where $e_t^\mathrm{price,full} = \sum_j [\Sigma_0^{-1}]_{1j}\sum_i\tilde x_t^{(i,1)}\tilde x_t^{(i,j)}$.

In angular terms: $\lambda_\mathrm{price}$ is the cosine of the angle between
the full energy gradient and the price-axis subspace of the Jacobian.
$\lambda_\mathrm{price} = 1$: all energy is in the price dimension — force aligned with price.
$\lambda_\mathrm{price} = 0$: energy in volatility/dispersion — force perpendicular to price.

This is the projection $\cos\angle(\gX, e_1)$ where $e_1$ is the price direction.

\begin{theorem}[Price loading as inner product]
\[
\lambda_\mathrm{price}(t) = \frac{|\inner{\gX(t)}{J_1(t)\T}|_F}{\norm{\gX(t)}_F\norm{J_1(t)\T}_F}
\]
where $J_1$ is the price-channel row of $J$.
\end{theorem}

\begin{proof}
$e_t^\mathrm{price,full}$ is the $(1,1)$ element of $\tilde X_t\Sigma_0^{-1}\tilde X_t\T$ restricted to the price-column Jacobian, which equals $\norm{J_1\T AS}^2 = \norm{\gX^{(C)}}^2/4$ in the pure-Mahalanobis case.
The ratio $\sqrt{e_t^\mathrm{price,full}/e_t}$ is therefore the Frobenius cosine between the price gradient component and the full gradient, as stated.
\end{proof}

%==============================================================
\part{Complete Angular Reference}
%==============================================================

\chapter{All Angles, All Formulas}

\section{The seven canonical angles}

\begin{resultbox}
\renewcommand{\arraystretch}{1.5}
\begin{tabular}{p{2.5cm}p{5.5cm}p{5cm}}
\toprule
\textbf{Angle} & \textbf{Formula} & \textbf{Role} \\
\midrule
Alignment $\theta_t$ & $\cos\theta_t = \inner{\gX}{\Fbase}/(\norm{\gX}\norm{\Fbase})$ & Energy descent condition \\
Misalignment $\psi_t$ & $\cos\psi_t = R_t/(\norm{\tilde G_t}\norm{g_t})$ & False alarm filter \\
Collapse $\tan\theta_t$ & $\tan\theta_t = \norm{F_\perp}/\norm{F_\parallel}$ & Collapse proximity \\
Angular velocity $\dot\theta_t$ & $-\dot h/\sin\theta_t$ (Theorem~\ref{thm:angvel}) & Rate of decoherence \\
Angular accel. $\ddot\theta_t$ & $(\theta_{t+1}-2\theta_t+\theta_{t-1})/\Delta t^2$ & Decoherence acceleration \\
Phase $\phi_i(t)$ & $\arg(V_1\T\tilde x_i + iV_2\T\tilde x_i)$ & Instantaneous agent phase \\
Approach angle $\alpha_\mathrm{app}$ & $\angle(\dot X, \nabla\kappa)$ & Critical manifold proximity \\
\bottomrule
\end{tabular}
\end{resultbox}

\section{Trigonometric identities in canonical form}

\begin{resultbox}
\begin{align*}
\norm{F_\parallel}^2 + \norm{F_\perp}^2 &= \norm{\Fbase}^2 \quad\text{(Pythagorean)} \\
\norm{F_\parallel} &= \norm{\Fbase}\cos\theta_t \quad\text{(projection)} \\
\norm{F_\perp} &= \norm{\Fbase}\sin\theta_t \quad\text{(rejection)} \\
\tan\theta_t &= \sin\theta_t/\cos\theta_t \quad\text{(collapse ratio)} \\
\cos\psi_t &= \mfls^\mathrm{state}/\mfls^\mathrm{channel} \quad\text{(MFLS ratio)} \\
\cos\theta_t &= \inner{\gX}{\Fbase}/(\mfls\norm{\Fbase}) \quad\text{(alignment)} \\
\dot\theta_t &= -\dot h/\sqrt{1-h^2} \quad\text{(angular velocity)} \\
\theta_t + \psi_t + \varphi_t &\leq \pi \quad\text{(angular triangle)} \\
E(R) &= 2\theta_\mathrm{rot}^2 \quad\text{(energy as squared angle on }SO(3)\text{)} \\
\gamma &= 1/(1+\theta/(2\theta_\mathrm{rot}^2)) \quad\text{(Fermi-Dirac in rotation angle)} \\
R(t) &= 1 - N^{-1}\textstyle\sum_i\delta_\Phi^{(i)} + O(1/N) \quad\text{(Kuramoto duality)} \\
\rho_\mfls &= \cos\psi_t \quad\text{(amplification factor)} \\
\lambda_\mathrm{price} &= \cos\angle(\gX, J_1\T) \quad\text{(price loading)}
\end{align*}
\end{resultbox}

\section{The three-phase precursor in angular notation}

\begin{lockedbox}
\textbf{Three-phase precursor} $\mathcal{P}(t)$:
\[
\underbrace{(1-\gamma)\left[\mfls\norm{\Fbase}\cos\theta_t - \mfls^2\right] > 0}_{\text{Phase 1: } \dot E > 0}
\;\wedge\;
\underbrace{\frac{d}{dt}\left|\frac{1}{N}\sum_i e^{i\phi_i}\right| > 0}_{\text{Phase 2: } \dot R > 0}
\;\wedge\;
\underbrace{\frac{\theta_{t+1}-2\theta_t+\theta_{t-1}}{\Delta t^2} > 0}_{\text{Phase 3: }\ddot\theta_t > 0}
\]
All three conditions stated in angle-space language.
$\mathcal{P}(t) \implies d\,\mathrm{SafetyRatio}/dt < 0$ (Theorem~\ref{thm:precursor}).
\end{lockedbox}

\section{Honest scope summary}

\begin{center}
\renewcommand{\arraystretch}{1.3}
\begin{tabular}{p{5.5cm}p{2.5cm}p{5cm}}
\toprule
\textbf{Claim} & \textbf{Status} & \textbf{Basis} \\
\midrule
$\theta_t$, $\psi_t$, $\tan\theta_t$ formulas & Proved & Direct from canonical definitions \\
Angular velocity exact formula & Proved & Chain rule \\
Phase-amplitude $\dot E$ split & Proved & Theorem~\ref{thm:phasplit} \\
Three-phase precursor & Proved & Theorem~\ref{thm:precursor} \\
Rodrigues formula & Proved & Standard exponential map \\
$\SO(3)$ energy as squared angle & Proved & Matrix logarithm identity \\
Kuramoto-coherence duality & Proved & Large-$N$ limit \\
Triple equivalence of collapse & Proved & Hessian + angular + energetic \\
Grid phase threshold & Proved & Hessian of Hamiltonian \\
Price loading as inner product & Proved & Direct computation \\
\midrule
$\ddot\theta_t > 0$ as trading signal & Testable & Computable from v58 history \\
Pre-collapse spectral shift & Plausible & Slow-mode argument (not rigorous) \\
Cardiac fibrillation prediction & Speculative & Framework applies, not validated \\
Turbulence from spectral phase & Speculative & Framework applies, not validated \\
\bottomrule
\end{tabular}
\end{center}

%==============================================================
\part{The Ellipsoidal Origin}
%==============================================================

\chapter{The Ellipsoid Was Already the Energy}

\section{The original ellipsoidal geometry}

The BSDT framework began with a critical manifold defined explicitly as an ellipsoid
in channel space. The stability boundary was the surface:
\[
\sum_k\left(\frac{\delta_k}{a_k}\right)^2 = 1,
\qquad a_k = \frac{\alpha}{w_k}
\]
where $\delta_k$ are the four BSDT channel scores and $w_k$ are Fisher variance
ratio weights.

This is a level set of the quadratic form $\delta^\top G\delta$ with metric:
\[
G = \operatorname{diag}\!\left(\frac{1}{a_C^2},\frac{1}{a_G^2},\frac{1}{a_A^2},\frac{1}{a_T^2}\right).
\]

The critical manifold was $\{\delta : \delta^\top G\delta = 1\}$.

\section{The canonical energy is the same object}

\begin{theorem}[Ellipsoid-Canonical identity]
\label{thm:origin}
The canonical energy $E = S^\top GS$, critical manifold $\mathcal{C}^* = \{E=1\}$,
and canonical gradient $\nabla_S E = 2GS$ are the energy-gradient reformulation
of the original ellipsoidal geometry. Specifically:
\[
\mathcal{C}^* = \{\delta : \delta^\top G\delta = 1\} = \{S : E(S) = 1\},
\]
\[
\rho := \sqrt{\delta^\top G\delta} = \sqrt{E},
\]
\[
n := G\delta = \tfrac{1}{2}\nabla_S E.
\]
The ellipsoid surface normal $n$ is exactly half the canonical gradient.
\end{theorem}

\begin{proof}
$E(S) = S^\top GS$. Setting $\delta = S$:
$E(\delta) = \delta^\top G\delta$,
so $\mathcal{C}^* = \{E=1\}$ and $\rho = \sqrt{E}$.
$\nabla_S E = 2GS = 2G\delta$,
so $n = G\delta = \frac{1}{2}\nabla_S E$.
$\square$
\end{proof}

\begin{remark}[What this means]
The canonical framework did not replace the ellipsoidal geometry.
It expressed the same geometry in gradient coordinates.
The ellipsoid $\sum_k(\delta_k/a_k)^2 \leq 1$ is the sublevel set $\{E \leq 1\}$.
Its surface is the level set $\{E = 1\}$.
Its outward normal is $\nabla_S E$.
Every subsequent canonical object is a consequence of this identification.
\end{remark}

\section{The surface-normal angle is the alignment angle}

In the early framework, the key diagnostic was the surface-normal angle:
\[
\cos\phi = \frac{\langle\dot\delta, n\rangle}{|\dot\delta||n|}
= \frac{\langle\dot S, \nabla_S E\rangle}{\|\dot S\|\|\nabla_S E\|}
\]
measuring how directly the trajectory moves toward the critical manifold.

\begin{theorem}[Surface-normal angle = alignment angle in channel space]
The surface-normal angle $\phi$ equals the canonical alignment angle $\theta_t$
evaluated in channel space:
\[
\cos\phi = \cos\theta_t\big|_\mathrm{channel} = \frac{\langle g_t, \Fbase\rangle}{\|g_t\|\|\Fbase\|}
\]
where $g_t = \nabla_S E_\mathrm{BS}(S_t)$ is the channel-space canonical gradient
and $\Fbase$ is identified with the channel-space velocity $\dot S$.
\end{theorem}

\begin{proof}
$\cos\phi = \langle\dot\delta, n\rangle/(|\dot\delta||n|)
= \langle\dot S, G S\rangle/(\|\dot S\|\|GS\|)
= \langle\dot S, \frac{1}{2}\nabla_S E\rangle/(\|\dot S\|\|\frac{1}{2}\nabla_S E\|)
= \langle\dot S, \nabla_S E\rangle/(\|\dot S\|\|\nabla_S E\|)
= \cos\theta_t|_\mathrm{channel}$.
\end{proof}

\section{Two states with the same $\rho$ but different futures}

This is the key insight from the ellipsoidal picture that the energy-only formulation
compressed away.

\begin{proposition}[Why $\phi$ carries information $\rho$ does not]
Two states $\delta_1, \delta_2$ with $\rho_1 = \rho_2$ (identical distance from centre,
identical energy $E$) can have:
\[
\cos\phi_1 = +1 \quad\text{(heading directly toward collapse)},
\]
\[
\cos\phi_2 = -1 \quad\text{(heading directly away from collapse)}.
\]
Their energy derivatives $\dot E = g(E)\langle\nabla E, \dot\delta\rangle$ have
opposite signs despite identical $E$. One is a genuine alarm; one is a recovery.
Energy alone cannot distinguish them. The angle $\phi$ does.
\end{proposition}

\begin{proof}
$\dot E = 2\langle G\delta, \dot\delta\rangle = \|\nabla E\|\|\dot\delta\|\cos\phi$.
Same $\rho$ implies same $E$ implies same $\|\nabla E\| = 2\|G\delta\|$.
Same $\|\dot\delta\|$ (equal velocity magnitudes).
$\cos\phi_1 = +1\Rightarrow\dot E_1 > 0$; $\cos\phi_2 = -1\Rightarrow\dot E_2 < 0$.
$\square$
\end{proof}

This explains why the alignment angle $\theta_t$ was added to the canonical dashboard:
it recovers the angular information that $E$ alone loses.

\section{The adaptive $\gamma$ origins}

The early GravityEngine used:
\[
\gamma^*_\kappa(X) \propto \frac{1}{\lambda_\mathrm{max}(\nabla^2\Phi(X))}
\]
— intervention proportional to landscape curvature.

The canonical framework replaced this with:
\[
\gamma_E(X) = \frac{E(X)}{E(X)+\theta}
\]
— intervention proportional to energy (distance from centre).

\begin{proposition}[Relationship between $\gamma_\kappa$ and $\gamma_E$]
At the critical manifold $\mathcal{C}^* = \{E=1\}$:
$\gamma_E = 1/(1+\theta)$.
At $\mathcal{C}_\mathrm{man} = \{\lambda_\mathrm{max}(\nabla^2\Phi)=1\}$:
$\gamma^*_\kappa = \alpha/(1+\varepsilon)$.
They coincide when $\theta = \varepsilon$ and both manifolds are the same set.

In general they diverge because $E$ measures displacement while $\kappa$ measures
landscape flatness. The landscape can flatten at moderate $E$
(agents approach the critical separation $r \approx \sigma$ in the erf-log potential),
causing $\kappa$ to spike while $E$ remains moderate.
This is why $\gamma^*_\kappa$ fires earlier: it detects the precondition
(flattening landscape) rather than the consequence (large displacement).
\end{proposition}

\section{The complete developmental chain}

\begin{lockedbox}
\textbf{The four-stage compression:}
\[
\underbrace{\text{Ellipsoid}}_{\delta^\top G\delta = 1}
\;\to\;
\underbrace{\text{Normal}}_{{n = G\delta}}
\;\to\;
\underbrace{\text{Gradient}}_{\nabla_S E = 2n}
\;\to\;
\underbrace{\text{Energy}}_{E = S^\top GS}
\]

\textbf{Exact correspondences:}
\begin{align*}
\mathcal{C}^* &= \{E = 1\} &&\text{(same set, two names)}\\
n &= \tfrac{1}{2}\nabla_S E &&\text{(same vector, factor of 2)}\\
\rho &= \sqrt{E} &&\text{(same scalar, square root)}\\
\cos\phi &= \cos\theta_t|_\mathrm{channel} &&\text{(same angle, different notation)}\\
\gamma^*_\kappa &\leftrightarrow \gamma_E &&\text{(same intent, Level 1 vs Level 3)}
\end{align*}

\textbf{What the compression preserved:}
descent structure, Lyapunov certificate, cross-domain generality.

\textbf{What the compression lost:}
explicit boundary orientation, curvature sensitivity, $\phi$ as a first-class object.

\textbf{What the canonical dashboard recovers:}
$\cos\theta_t$ recovers $\phi$;
$\kappa_\mathrm{norm}$ recovers $\gamma^*_\kappa$;
$\psi_t$ is genuinely new (no predecessor in the ellipsoidal framework).
\end{lockedbox}

\section{The dual description theorem}

\begin{theorem}[Ellipsoid-canonical duality]
The ellipsoidal BSDT geometry and the canonical energy framework are
dual descriptions of the same quadratic geometry.

The ellipsoidal description is boundary-oriented:
it monitors distance to $\mathcal{C}^*$ and direction of approach.

The canonical description is centre-oriented:
it monitors distance from the normal-period mean and the energy growth rate.

Both are equivalent at the critical manifold. They differ in what they emphasise
away from the manifold:
the ellipsoidal framework is more sensitive to approach geometry at moderate $\rho$;
the canonical framework is more tractable for Lyapunov proofs and cross-domain generalisation.

The complete framework uses both: canonical energy for control,
ellipsoidal boundary orientation ($\cos\theta_t = \cos\phi$) for diagnosis.
\end{theorem}

\begin{proof}
Theorem~\ref{thm:origin} establishes the exact identifications.
The dual description follows: same set $\mathcal{C}^*$, same normal $n \propto \nabla E$,
same angle $\phi = \theta_t|_\mathrm{channel}$.
The distance emphasis difference ($1-\rho$ vs $E$) is a coordinate choice,
not a structural difference.
$\square$
\end{proof}

\section{Summary: what was always there}

Every object in the canonical framework has an exact predecessor in the
ellipsoidal geometry:

\begin{center}
\renewcommand{\arraystretch}{1.4}
\begin{tabular}{lll}
\toprule
\textbf{Ellipsoidal framework} & \textbf{Canonical framework} & \textbf{Relation} \\
\midrule
$\mathcal{C}^* = \{\delta^\top G\delta=1\}$ & $\{E=1\}$ & Same set \\
$n = G\delta$ (surface normal) & $\frac{1}{2}\nabla_S E$ & $n = \frac{1}{2}\nabla_S E$ \\
$\rho = \sqrt{\delta^\top G\delta}$ & $\sqrt{E}$ & $\rho = \sqrt{E}$ \\
$\cos\phi = \langle\dot\delta,n\rangle/(|\dot\delta||n|)$ & $\cos\theta_t|_\mathrm{ch}$ & Identical \\
$\gamma^*_\kappa \propto 1/\lambda_\mathrm{max}$ & $\gamma_E = E/(E+\theta)$ & Level 1 vs Level 3 \\
$\psi$ (boundary approach rate) & $\dot E = g(E)\langle\nabla E,\Fbase\rangle$ & $\dot E = \|\nabla E\|\|\Fbase\|\cos\phi \cdot g(E)$ \\
\midrule
\textit{No predecessor} & $\psi_t$ (misalignment angle) & Genuinely new \\
\textit{No predecessor} & $R(t)$ (Kuramoto order parameter) & Genuinely new \\
\textit{No predecessor} & $\ddot\theta_t$ (angular acceleration) & Genuinely new \\
\bottomrule
\end{tabular}
\end{center}

The ellipsoidal framework contained the energy, the gradient, the critical manifold,
and the alignment angle --- all implicitly.
The canonical framework made them explicit, proved stability, and generalised
across domains.
The three genuinely new objects ($\psi_t$, $R(t)$, $\ddot\theta_t$) have no
ellipsoidal predecessors: they emerged from the canonical abstraction itself.

\end{document}
